\documentclass[11pt]{article}
\newcommand{\ReviewCode}{DRAFT}
% Shared typesetting; \\input from a project two directories under the root.
\usepackage[T1]{fontenc}
\usepackage[utf8]{inputenc}
\usepackage{lmodern}
\usepackage[letterpaper,margin=1.12in,headheight=14pt]{geometry}
\usepackage{microtype}
\usepackage{amsmath,amssymb,amsthm,mathtools,bm}
\usepackage{booktabs,array,longtable}
\usepackage{graphicx}
\usepackage{xcolor}
\usepackage{enumitem}
\usepackage{csquotes}
\definecolor{ink}{HTML}{172535}
\definecolor{linkblue}{HTML}{1B527D}
\definecolor{muted}{HTML}{536271}
\usepackage[backend=biber,style=authoryear,sorting=nyt,maxbibnames=99,doi=true,url=true,isbn=false,backref=true]{biblatex}
\addbibresource{../../bibliography/references.bib}
\usepackage[pdfusetitle,colorlinks=true,linkcolor=ink,citecolor=linkblue,urlcolor=linkblue]{hyperref}
\usepackage{bookmark}
\usepackage[nameinlink,noabbrev]{cleveref}
\usepackage{fancyhdr}
\pagestyle{fancy}
\fancyhf{}
\fancyhead[L]{\small\textsf{RESEARCH COMPENDIUM}}
\fancyhead[R]{\small\textsf{\ReviewCode}}
\fancyfoot[C]{\thepage}
\renewcommand{\headrulewidth}{0.3pt}
\setlength{\parindent}{0pt}
\setlength{\parskip}{5pt}
\setcounter{tocdepth}{2}
\newtheorem{definition}{Definition}[section]
\newtheorem{remark}{Remark}[section]
\newcommand{\ReviewStatusText}[1]{\par\noindent\textcolor{muted}{\small\textsf{#1}}\par\medskip}
% Unambiguous, reusable physical notation.
\newcommand{\dd}{\mathop{}\!\mathrm{d}}
\newcommand{\ii}{\mathrm{i}}
\newcommand{\ee}{\mathrm{e}}
\newcommand{\ket}[1]{\lvert #1\rangle}
\newcommand{\bra}[1]{\langle #1\rvert}
\newcommand{\braket}[2]{\langle #1\vert #2\rangle}
\newcommand{\abs}[1]{\left\lvert #1\right\rvert}
\newcommand{\norm}[1]{\left\lVert #1\right\rVert}
\newcommand{\op}[1]{\widehat{#1}}
\DeclareMathOperator{\ImPart}{Im}
\DeclareMathOperator{\RePart}{Re}


\title{\bfseries Review Title\\[3pt]\large A Precise Research Question}
\author{Research Compendium \;\textbar\; Working review}
\date{\today}
\begin{document}
\maketitle
\ReviewStatusText{STATUS: OUTLINE \quad\textbullet\quad Do not cite as a completed review}
\begin{abstract}
State the physical question, the set of admissible observations, the hypothesis space, and the limits of the present review. Avoid unsupported quantitative claims.
\end{abstract}
\tableofcontents
\clearpage
\section{Question and scope}\label{sec:question}
Define the physical problem and the exact quantities to be inferred.
\section{Conventions}\label{sec:conventions}
Define kinematics, normalization, and signs before using results.
\section{Evidence}\label{sec:evidence}
List experimentally observed quantities and their sources, distinguishing direct from secondary analysis.
\section{Models and interpretation}\label{sec:theory}
Record dynamical assumptions and what the model can or cannot infer.
\section{Open problems}\label{sec:open}
Identify a testable next measurement or calculation.
\section*{Revision log}
\addcontentsline{toc}{section}{Revision log}
\textbf{\today}: initial outline.
% Cite authoritative literature with \textcite{verified-key} as you write.
\printbibliography[heading=bibintoc,title={References}]
\end{document}
